\cotefile{137}
\scanpage{2}
$k$ anneau commutatif, $k_*$ \unsure{le $k$ simplicial|le demi-groupe|est l'anneau $\frac12$ simpl.} \unsure{défini}.

\uline{Lemme 1} Soit $L_*$ un $k$-module $\frac12$ simplicial, \unsure{plat et f.g.}

Pour que $L_*$ soit un $k_*$-Module projectif, il suffit qu'il soit flasque (i.e. $\pi_i(L_*)=0$ $\forall i\geq0$) et que les $L_n$ soient des $k$-modules projectifs.
% margin note written sideways
\margin{\unsure{p.ex. il} suffit \unsure{que les} $\unsure{L_q(L_*)}$ \unsure{soient plats}.}

La \unsure{condition|complexe} signifie aussi que le complexe des chaînes (des $k$-modules) associé par Dold-Puppe est acyclique et à \unsure{compos.} projectives, on voit alors qu'il \deleted{est} \inserted{\unsure{devient}} isomorphe à une somme directe de complexes de chaînes de la forme
\[
\ldots\,0\longrightarrow0\longrightarrow L_q
 \xrightarrow{\unsure{\sim}}L_q\longrightarrow0\longrightarrow0\,\ldots
\]
\unsure{non dégénérés|non dégénéres} \unsure{à fois conséc.|en degrés conséc.}, en $\unsure{k}$. Or un tel complexe est évidemment projectif (dans la cat. des complexes de chaînes de $k$-modules), \unsure{donc aussi une somme quelconque de tels}, d'où on conclut \unsure{par l'équivalence de catégories}.

\uline{NB} Il suffit, au lieu de supposer $L_*$ flasque i.e. $\pi_i(L_*)=0$ $\forall i\geq0$, de supposer $\pi_i(L_*)=0$ pour $i>0$ et $\pi_0(L_*)$ $k$-projectif. (On trouve \unsure{des sommes directes des complexes|alors sans doute une condition} \unsure{en} \unsure{ci et s.}) En effet, dans l'argument précédent, il faut de plus tenir compte \unsure{d'un sommande|des sommandes} qui est un module projectif $L_0$ placé en degré $0$. Or il est immédiat que c'est encore un complexe de chaînes projectif.

\uline{Cor.} Les $\Gamma^i\Phi_{*,k}\otimes_k\Lambda^j\Psi_{*,k}$ \inserted{et $\Lambda^i\Phi_{*,k}\otimes\Lambda^j\Psi_{*,k}$} (pour $i>0$ sont des modules $\frac12$ simpliciaux projectifs.
% archive sheet number: 1

\scanpage{3}
\mspage{112.}
\uline{APPENDIX} \uline{H} (\uline{Hirsch Method})

If $G\to E\to M$ is a principal $G$-bundle over a manifold $M$, and $\handwritten{\mathscr A}$ is a differential algebra mapping to the forms on $M$ which computes the cohomology of $M$, then
\[
\handwritten{\mathscr A}\otimes\Lambda(x_1,\ldots,x_n)
\]
can be supplied with a differential in such a way that the cohomology of $E$ results. Here $\Lambda(x_1,\ldots,x_n)$ is the cohomology of the fibre $G$ and $d(1\otimes x_i)$ is $c_i\otimes1$ where $c_i$ is a representative of the characteristic class of $E$.

This same method is applied inductively to Postnikov systems to prove the theorems of the paper.

Now we are inductively considering principal fibrations
\diagram{\begin{tikzcd}[ampersand replacement=\&]
K\arrow[r] \& X_{n+1}\arrow[d,"{\deleted{\unsure{\pi|n}}\handwritten{\pi_n}}"]\\
\& X_n
\end{tikzcd}}
where $K$ has one non-zero Abelian homotopy group in dimension $i_n\geq1$.

Using the Hirsch method and the existence of the rational de Rham algebra which is commutative and correctly computes cohomology we can inductively construct differential algebras of the minimal model type which compute the cohomology of the spaces $\{X_{n+1}\}$ in any countable system of maps like the $\pi_n$.

At each stage in the construction we are adding a vector space of generators of the same rank and dimension of the homotopy group of $K$ to the algebra for $X_n$ to obtain
% archive sheet number: 2

\scanpage{4}
\uline{Lemme 2} Soit $L_*$ un $k$-module $\frac12$ simplicial, \deleted{\unsure{tel que}} \deleted{\unsure{(projectif)}} \deleted{flasque,} \deleted{que,} $q\geq q'\geq0$ deux entiers, $M$ un $k$-module, \inserted{\unsure{on suppose}, si $q=q'$, $\pi_{\unsure{0}}(L_*)=0$.} Alors
\[
\operatorname{Hom}_{\unsure{\widehat\Delta\mathrm{Ab}|\mathrm{Zab}}}
 (\underbrace{L_*\otimes_k\Lambda^q\Psi_{*,k}}_{\substack{\text{(on utilise seult}\\\text{que ce complexe est nul en dim }<q\text{)}}},\Lambda^{q'}\Psi_{*,k}\otimes_k M)=0.
\]


\deleted{Si $L'_*$ est le \unsure{complexe de chaînes associé}} \inserted{Le complexe de chaînes associé à} \deleted{\unsure{$L_*$, le complexe associé à}} $L_*\otimes_k\Lambda^q\Psi_{*,k}$ est \deleted{\unsure{un complexe}} \deleted{\unsure{de $k$-modules gradués, donc}} de la forme
\[
\ldots\longrightarrow L'_{q+1}\longrightarrow L'_q\longrightarrow0\longrightarrow0\ldots
\]
tandis que $\Lambda^{q'}\Psi_{*,k}\otimes_k M$ est $M$ placé en degré $q'$, \inserted{i.e. $M[q']$} si $q'<q$, il est évident que $\operatorname{Hom}_{\unsure{\bullet}}(L'_{\bullet},M_{\bullet})=0$, pour des raisons de \unsure{degrés|degré}. Si $q'=q$,
\diagram{\begin{tikzcd}[ampersand replacement=\&]
L'_{q+1}\arrow[r,"{\mathrm{surj}}"]\arrow[d] \& L'_q\arrow[r]\arrow[d] \&0\,\ldots\\
0\arrow[r]\& M\arrow[r]\&0
\end{tikzcd}}
alors $\pi_{\unsure{0|q}}(L_*)=0\ \Longrightarrow\ L'_{q+1}\xrightarrow{d}L'_q$ est surjectif, donc $L'_q\longrightarrow M$ est \unsure{nul|nulle} pour tout hom. de complexes \inserted{de chaînes} $L'_{q}\longrightarrow M[q]$, cqfd.

\uline{Théorème} Soit $n\geq0$. Il existe un homomorphisme de complexes de cochaînes dans $\unsure{\widehat{\Delta}_{\mathrm{ab}}}$
\diagram{\begin{tikzcd}[ampersand replacement=\&,column sep=tiny,font=\scriptsize]
0\arrow[r]\& \overset{0}{\Gamma^n\Phi_*}\arrow[r]\arrow[d,"u_0"]\&
\Gamma^{n-1}\Phi_*\otimes\Psi_*\arrow[r]\arrow[d,"u_1"]\&
\Gamma^{n-2}\Phi_*\otimes\Lambda^2\Psi_*\arrow[r]\arrow[d,"u_2"]\&
\cdots\arrow[r]\& \Lambda^n\Psi_*\arrow[r]\arrow[d,"u_n"]\&0\arrow[r]\arrow[d]\&0\\
0\arrow[r]\&\Phi_*\arrow[r,"{\wedge T}"']\&\Lambda^2\Phi_*\arrow[r,"{\wedge T}"']\&
\Lambda^3\Phi_*\arrow[r]\&\cdots\arrow[r,"{\wedge T}"']\&
\Lambda^{n+1}\Phi_*\arrow[r,"{\wedge T}"']\&\Lambda^{n+2}\Phi_*\arrow[r]\&\cdots
\end{tikzcd}}
(nb : on écrit $\Phi_*$ et $\Psi_*$ pour $\Phi_{*k}$, $\Psi_{*k}$, et $\otimes$ signifie $\otimes_k$), tel que $u_n$ soit l'isom. can. de $\Lambda^n\Psi_*$ sur $\operatorname{Ker}(\Lambda^{n+1}\Phi_*\twoheadrightarrow\Lambda^{n+1}\Psi_*\hookrightarrow\Lambda^{n+2}\Phi_*)=_{\mathrm{df}}Z(\Lambda^{\unsure{n+1}}\Phi_*)$
% archive sheet number: 3

\scanpage{5}
The map $\phi$ of $M$ to the compact nil-manifold $H/\Gamma$ also generalizes and we obtain a tower of maps of $M$ into a tower of nil-manifolds beginning with the Jacobian $H_1(M,R)/H_1(M,Z)$.

The next section shows how this construction becomes canonical if $M$ is Riemannian

\uline{The} \uline{Model} \uline{of} \uline{a} \uline{Riemannian} \uline{Manifold}

If $M$ is Riemannian we have the $*$-operator on forms, the adjoint of $d$, $d^*=*d*$, and the Hodge decomposition
\[
\omega=dx_\omega+d^*y_\omega+h_\omega
\]
where $x_\omega$, $y_\omega$, and $h_\omega$ are unique subject to the conditions $x_\omega\unsure{\varepsilon|\in}\text{ image of }d^*$, $y_\omega\varepsilon\text{ image }d$, $h_\omega\varepsilon\ker d\handwritten{\cap}\ker d^*$.

This follows on the formal level from the \textquotedbl{}disjointness\textquotedbl{} of $d$ and $d^*$:
\[
\begin{aligned}
dd^*x=0&\unsure{\Rightarrow}d^*x=0\\
d^*d\,y=0&\Rightarrow dy=0.
\end{aligned}
\]

Thus $\omega$ is closed iff $y_\omega=0$. Then $\omega$ is cohomologous to a canonical form, the harmonic representative $h_\omega$. And finally if $\omega$ is exact $y_\omega$ and $h_\omega$ are zero and $\omega$ is the differential of a canonical form $x_\omega$.

We can use the Hodge decomposition to rigidify the model of a Riemannian manifold. In the construction of $\rho$
\[
\Lambda(x_1,x_2,\ldots;d)\xrightarrow{\rho}\{\text{smooth forms}\}
\]
we have to make various choices. For example we have to choose
% archive sheet number: 4

\scanpage{7}
\uline{Lemme} Supposons les $L_n$ et $\pi_*(L_{\bullet})$ projectifs de t.f. sur $k$, et supposons $k$ commutatif à \uline{spec connexe} (donc $\forall M$ projectif sur $k$, $\operatorname{rg}_k M\in\mathbb N$ est défini).
% slanted
\margin{on peut \unsure{évid$^{\mathrm{mt}}$} \unsure{éliminer cette|changer cette|éviter cette} hypothèse}

Conditions équivalentes

a) $L_{\bullet}^{!}$ à degrés bornés

b) \deleted{\unsure{$L_{\bullet}$|Tous}} Les $M_i,N_j$ \unsure{n.} nuls sont en nb fini i.e. $L_{\bullet}$ est \inserted{$\simeq$} somme \uline{finie} de modules $\frac12$ simpliciaux $\Lambda^i\Phi_*\otimes M_i$ $(i\geq1)$ $\Lambda^j\Psi_*\otimes N_j$ $(j\geq1)$ ($M_i$ et $N_j$ proj. de t.f.)

b') (Si $k$ principal, plus gén si tout module proj de t.f. sur $k$ est libre). $L_{\bullet}$ est \inserted{$\simeq$} somme \uline{finie} de modules $\frac12$ simpliciaux $\Lambda^i\Phi_*$ \inserted{$(i\geq1)$}, $\Lambda^j\Psi_*$ $(j\geq0)$

c) $\operatorname{rang}_k L_n$ est une fonction polynômiale en $n$ pour $n$ grand

d) $\exists N$ \inserted{et $c>0$} tel que pour $n$ ass. grand, on ait
\[
\operatorname{rang}_k L_n\leq c\,n^N
\]

\deleted{\unsure{Cet}} [De plus, \deleted{p} pour $N$ fixé, la condition que l'on ait $\operatorname{rang}_k L_n=O(n^N)$ signifie que $M_i,N_j$ sont nuls pour $i,j>N$ ; \deleted{et si $N$ est choisi le plus petit possible (i.e. tel que $M_N+N_N\ne0$), \unsure{la limite est un multiple de $1/N!$}.})
% slanted
\margin{et on a plus $\displaystyle\lim_{n\to\infty}\frac{\operatorname{rg}_k L_n}{n^N}=\mu_N+\nu_N$}

\uline{Dém} \deleted{Equivalence de a) et b)} $a)\Longleftrightarrow b)$ déjà vu (voir 1 et 2 ci-dessus), $b)\Longleftrightarrow b')$ trivial, \deleted{$a)\Longrightarrow$} $c)\Longrightarrow d)$ trivial, $a)\Longrightarrow c)$ et $d)\Longrightarrow a)$ immédiat par la formule explicite
% archive sheet number: 6

\scanpage{8}
\uline{Prop} Soit $L_{\bullet}$ un $k$-module $\frac12$ simplicial. Les conditions suivantes sont équivalentes

a) les $L_n$ et $\pi_n(L_{\bullet})$ sont des $k$-modules projectifs.

\deleted{b}a'{)} Les composantes \deleted{de} $(L_{\bullet}^!)_n$ de $L_{\bullet}^!$ (le complexe de chaînes de Dold-Puppe associé à $L_{\bullet}$) et les $H_n(L_{\bullet}^!)$ sont des $k$-modules projectifs

b) $L_{\unsure{*|\bullet}}$ est isomorphe à \unsure{une|un} \deleted{somme} \inserted{produit} direct, \deleted{de modules} de la forme
\[
\left(\prod_{i\geq1}\Lambda^i\Phi_{*k}\otimes_k M_i\deleted{\unsure{d}}\right)
 \times\prod_{j\geq\deleted{\unsure{1}}0}(\Lambda^j\Psi_{*k})\otimes_k N_j
\]

\deleted{\unsure{d|a}} b') $L^!$ est isomorphe à un produit direct (de complexes de cochaînes) de la forme
\[
\prod_{i\geq1}(0\longrightarrow\underbrace{M_i
 \unsure{\mathrel{-}|\longrightarrow}M_i}_{\text{en degré }i\ \text{et }i-1}
 \longrightarrow0\longrightarrow\ldots)
 \times\prod_{j\geq0}(0\longrightarrow
 \underset{\text{en degré }j}{N_j}\longrightarrow0\longrightarrow\ldots)
\]

De plus, les $M_i$ et $N_j$ de b) (resp b') sont déterminés à isom près par \deleted{\unsure{$K$}} $L_{\bullet}$, et \deleted{ce} sont les $\overline{\mathrm{m}}$\unsure{s} (à isom près) dans b) et b').

\uline{Dém} (\deleted{\unsure{b}}a'{)}$\Longleftrightarrow$b')) est trivial et bien connu, ainsi que l'unicité à isom \inserted{can} près des $M_i,N_j$ dans b') (\uline{NB} on a $N_j\simeq H_j(L_{\bullet})$, et $M_i\simeq L_i/Z_i\{\simeq B_{i-1})\deleted{\unsure{\simeq}}$. On a a)$\Longleftrightarrow$a') \inserted{et b)$\Longleftrightarrow$b')} par le \unsure{tapis} de Dold-Puppe. D'où l'équivalence des 4 conditions, et la dernière assertion est contenue.
% archive sheet number: 7

\scanpage{9}
% Pre-ruled music paper.
\uline{Cor} 1 \deleted{\unsure{Il en sera}} \deleted{\unsure{Dans ces conditions}} On obtient encore des conditions équivalentes en exigeant \inserted{de plus} sous a) que les $L_n$, \deleted{$\pi_n(L_{\bullet})$} soient de type fini, en a') que les $L_n^!$, \deleted{$H_n(L_{\bullet}^!)$} soient de type fini, dans b) et b') que les $M_i$ et $N_j$ soient de type fini.

On \unsure{aura|a alors} dans les conditions de prop 1
\[
 L_n=\prod_{1\leq i\leq n+1}\Lambda^{i}\Phi_{n k}\otimes_k M_i
 \times\prod_{\unsure{0|1}\leq j\leq n\deleted{\unsure{+1}}}\Lambda^j\Psi_n\otimes N_j
\]
\deleted{d'où on trouve}

\uline{Cor} 2 Sous les conditions équivalentes \deleted{de} \deleted{la} \deleted{\unsure{la prop., pour que $L_n$ soit de type fini, pour}} \inserted{\unsure{précédentes}, pour que $L_{\bullet}^!$ soit} à degrés bornés, il faut\unsure{)}, qu'il n'y ait qu'un nombre \inserted{fini} d'indices $i$ et $j$ tels que $M_i\ne0$, $N_j\ne0$.

\deleted{Sous les conditions} On appellera \unsure{les|un} $L_{\bullet}$ \deleted{\unsure{strict}} \deleted{\unsure{parfait}} \deleted{\unsure{spécial}} \inserted{\uline{\unsure{spécial}}} si $L_{\bullet}^!$ est \deleted{str. parfait, i.e.} \inserted{à degrés bornés, à composantes} et à $H_n$ projectifs de t.f. --- i.e. si $L_{\bullet}$ est \inserted{$\simeq$} somme finie de modules de la forme \deleted{\unsure{$k$}} $\Lambda^{\unsure{j|i}}\Phi_{*k}\otimes_k M_i$ \inserted{$(i\geq1)$}, $\Lambda^j\Psi_*\otimes M_j$ $(j\geq0)$, avec les $M_i$ et $N_j$ proj. de t.f., \uline{strictement spéciaux} si de plus les $M_i,N_j$ sont \uline{libres}. Il revient au $\overline{\mathrm m}$ de
% archive sheet number: 8

\scanpage{10}
% Pre-ruled music paper.
dire que $L_{\bullet}\simeq L_{\bullet}^{\circ}\otimes_{\mathbb Z}k$, où $L_{\bullet}^{\circ}$ est un complexe de $\mathbb Z$-modules $\frac12$ simplicial spécial.

\uline{Proposition 2} Les \deleted{\unsure{$\mathbb Z$}}$k$-modules \inserted{$\frac12$ simpliciaux} spéciaux (resp. str. spéciaux) sont stables par produits tensoriels [et par les op. tensorielles $\Lambda^i,\Gamma^i,\operatorname{Sym}^i$ ?]
% slanted
\margin{\unsure{lemme}}
% slanted
\margin{[ ] faux}

Le cas strictm spécial se ramène au cas spécial, en se plaçant sur $\mathbb Z$\ldots

1°) Produit tensoriel. On est réduit à prouver que les modules
\[
\left\{\begin{aligned}
\Lambda^i\Phi_*^{\deleted{\unsure{*}}}\otimes\Lambda^j\Phi_* &\qquad(i,j\geq1)\\
\Lambda^i\Psi_*\otimes\Lambda^j\Psi_* &\qquad(i,j\geq0)\\
\Lambda^i\Phi_*\otimes\Lambda^j\Psi_* &\qquad(i\deleted{,j}\geq1,j\geq0)
\end{aligned}\right.
\]
sont spéciaux. Or ils sont à composantes libres de t.f., et par Künneth-Eilenberg-Zilber leurs $\pi_i$ sont itou, il reste à montrer qu'il n'y a qu'un nb fini de facteurs $M_i,N_j\ne0$. \deleted{On va} \inserted{Cela résulte} \deleted{prouver} aussitôt du

\deleted{\uline{Lemme} \deleted{Sous} Les conditions du cor. 2 de prop 1 sont encore équivalentes à \unsure{celles|les|la} suivantes (en supposant \inserted{$k$ comm. et} $\operatorname{Spec}k$ connexe, de sorte que $\mu_i=\operatorname{rg}M_i$ et $\nu_j=\operatorname{rg}N_j$ est défini)

a) $\operatorname{rg}_k L_n$ est un polynôme en $n$ pour $n$ grand

b) $\exists$}
% archive sheet number: 9

\scanpage{11}
\[
\rho(n)\overset{\mathrm{dfn}}{=}\operatorname{rg}_k L_n
=\sum_{1\leq i\leq n+1}\mu_i
\underset{\substack{\downarrow\\\text{pol. de degré}\\i\ \text{en }n}}{\binom{n+1}{i}}
+\sum_{0\leq j\leq n}\nu_j
\underset{\substack{\downarrow\\\text{pol. de degré}\\j\ \text{en }n}}{\binom n j}
\]
où $\mu_i=\operatorname{rg}_k M_i$, $\nu_j=\operatorname{rg}_k N_j$

donc si \deleted{$\mu_i$ (resp $\nu_j$) \unsure{$=$}} un $\mu_i$ ou un $\nu_j$ est $\ne0$, cela montre que $\rho(n)$ pour $n$ grand est minoré par un polynôme de degré $i$ resp $j$.

[2°) Cas de $\Lambda^k$. Utilisant la formule de $\Lambda^k$ d'une somme directe, et 1°), on est ramené à prouver que les $\Lambda^k$ d'un module $\frac12$ simplicial $L_{\bullet}$ de la forme $\Lambda^i\Phi_*\otimes M_{\unsure{i|{}}}$, $\Lambda^j\Psi_*\otimes N$. Comme $M,N$ est \unsure{facteur|facteurs} direct d'un libre $P$, $L_{\bullet}$ est facteur direct de $L'_{\bullet}=\Lambda^i\Phi_*\otimes P$ ou $\Lambda^j\Psi_*\otimes P$, \deleted{\unsure{et que}} \deleted{\unsure{prop. 1 ci-dessus est}} et en réutilisant la formule de $\Lambda^k$ d'une somme directe, \unsure{montrant|montre} que $\unsure{\Gamma|\Lambda}^k(L_{\bullet})$ est facteur direct de $\Lambda^k L'_{\bullet}$. \inserted{Comme la propriété du \unsure{lemme} \deleted{\unsure{de}} de \unsure{prop.} 1, son cor. 1 et du cor 2 \deleted{sont} stable par fact directs,} on est ramené à prouver $\Lambda^i L'_{\bullet}$ est \deleted{\unsure{str}} spécial, donc \unsure{sps|supp.} $M,N$ libres --- donc $\overline{\mathrm m}$ (par la formule de la somme) $M=k$, $N=k$. Donc on est ramené à prouver
% archive sheet number: 10

\scanpage{12}
% Pre-ruled music paper.
$\Lambda^k(\Lambda^i\Phi_{*k})$ et $\Lambda^k(\Lambda^j\Psi_{*k})$ sont spéciaux (bien sûr, ils seront $\overline{\mathrm m}$ \unsure{strts|très} spéciaux).

Leurs composantes sont évidemment libres. De plus, comme $\Lambda^i\Phi_{*k}$ est homotope à $0$, de $\overline{\mathrm m}$ $\Lambda^k(\Lambda^i\Phi_*)$, dont les $\pi_i$ sont donc nuls et à fortiori libres de t.f. :
\margin{?}
\hfill $]$

\uline{Remarque}. Soit $L_{\bullet}$ $k$-module $\frac12$ simplicial \deleted{\unsure{str}} spécial, alors $L_{\bullet}$ est connu à isom près, i.e les $\mu_i=\operatorname{rg}M_i$ $(i\geq1)$ et $\nu_j=\operatorname{rg}N_j$ \inserted{$j\geq0$,} sont connus, quand on connait la \deleted{\unsure{valeurs}} \inserted{fonction} $n\longmapsto\operatorname{rg}_k L_n=f_{L_{\bullet}}(n)$, \deleted{En effet} \uline{et} les $\operatorname{rg}_k\deleted{\unsure{H}}\pi_j(L_{\bullet})$ $(=\operatorname{rg}_k N_j=\nu_j)$. En effet, on aura
\[
\sum_{i\geq1}\mu_i\underset{\substack{\Vert\\P_i(n+1)}}{\binom{n+1}i}
=f_{L_{\bullet}}(n)-\sum_{j\geq0}\nu_j\binom n j
\qquad(n\ \text{grand})
\]
\deleted{Or on sait} i.e.
\[
f_{L_{\bullet}}(n-1)-\sum_{j\geq0}P_j(n-1\deleted{\unsure{j}})=\sum_{i\geq1}\mu_i P_i(n)
\]
or on sait que les \inserted{pol $P_i(n)$} \deleted{\unsure{$\binom n i$}} sont \deleted{des} lin. \unsure{indépendts|indépendants} --- et $\overline{\mathrm m}$ forment une base du $\mathbb Z$-module des polynômes à \inserted{coeff. entiers et} valeurs entières sur les entiers.

Donc la connaissance des $\sum_{i\geq1}\mu_iP_i$ pour $n$ grand implique
% archive sheet number: 11

\scanpage{13}
la connaissance des $\mu_i$. Si on se donne les \inserted{entiers} $\nu_j\geq\unsure{0}$ \deleted{$\unsure{\mathrm{rg}\,\pi_j(L_\bullet)|=\mathrm{rg}\,\pi_j(L_\bullet)}$} \deleted{$\unsure{(j\geq0)}$} $(j\geq0)$ et une fonction \inserted{polynômiale} $f(n)$ à valeurs entières définie \inserted{sur les entiers} \deleted{pour $n$ grand}, à quelles conditions existe-t-il $L_\bullet$ très spécial avec $\nu_j=\mathrm{rg}\,\pi_j(L_\bullet)$ $(j\geq0)$ et $f(n)=f_{L_\bullet}(n)$ pour $n$ grand ? Notons que $g=f-\sum_{j\geq0}\nu_jP_j$ est encore une fonction \deleted{\unsure{entière}} \inserted{polynômiale} à valeurs entières, donc de la forme $\sum_{i\geq0}\mu_i\binom{n+1}{i}$ où $(\mu_i)\in\mathbb Z^{(\mathbb N)}$ est uniquement déterminé. Alors la condition demandée est
\[
\mu_0=0,\quad\mu_i\geq0\quad\text{si }i\geq1.
\]
Notons que \deleted{$\unsure{\mu_0=|\mu_0}$}
\[
\mu_0=g(-1)=f(-1)-\sum_{j\geq0}\nu_j\frac{(-1)(-2)\cdots(-j)}{j!}
\]
donc on a \unsure{bien la|encore}
\[
(\mu_0=0)\Longleftrightarrow\left(f(-1)=\underset{\substack{\Vert\\\chi((\nu_j))}}{\sum(-1)^j\nu_j}\right)
\]

\deleted{Déf. \deleted{Complexes} $k$-Modules $\frac12$ simpliciaux $L_\bullet$ \deleted{\unsure{prép.}} \inserted{\unsure{str.}} ``chaîne-parfait'' \inserted{(strictly chain-perfect)} $={}$ tels que $L_\bullet^!$ soit un complexe de chaînes de $k$-modules \deleted{\unsure{str.}} parfait (str. parfait) $={}$}

% archive sheet number: 12

\scanpage{14}
$k$ anneau (com. unit.)

\uline{Déf 1} Un $k$-module $\frac12$ simplicial \inserted{$L_\bullet$} est dit $\partial$-pseudo-cohérent (resp. strictement $\partial$-pseudocohérent, resp. $\partial$-parfait, resp strictement $\partial$parfait) si $L_\bullet^!$ (le complexe de chaînes associé par Dold-Puppe) est pseudo-cohérent (resp. str. ps. cohérent, resp parfait, resp. strictt parfait). Donc \unsure{de-par|de par} déf a) \inserted{ci-dessous}

\uline{Prop 1} a) $L_\bullet$ est str. $\deleted{\unsure{\phi}}\partial$-pseudocohérent (resp. str. $\partial$-parfait) si et seulmt si les $(L_n^!)_{\unsure{n}}$ sont projectifs de t. f. (resp. et nuls sauf un nb fini), \deleted{\illegible{2}}

b) $L_\bullet$ est str. $\partial$-ps. coh. \unsure{sss|ssi} les $L_n$ sont projectifs de t. f.

c) $L_\bullet$ est str. $\partial$parfait sss $L_\bullet$ est extension successive \inserted{finie} de Modules de la forme $\Lambda^i\Psi_\bullet\otimes_{\mathbb Z}M_i$ \inserted{$(i\geq0)$}, où les $M_i$ sont $k$-modules pr. de t.f. (l'ordre des facteurs dans une filtration croissante \unsure{peut être pris|peuvent être|pouvant} d'ailleurs être choisi par ordre des $i$ croissants), ou encore \deleted{\unsure{sss|ssi}} \inserted{(ce qui suit se \unsure{comprend} cas spécial : commencer par s'y placer)}
% slanted
\margin{sss les $L_n$ sont proj. de t. f. et $f(n)=\mathrm{rg}_k L_n$ est une fonction polynômiale pour $n$ grand (on montrera \unsure{ici} une telle fonction)}

d) Si $k$ est noeth., $L_\bullet$ est $\partial$-pseudo-cohérent (sss les $L_n^!$ sont de type fini sur $k$, ou encore) sss les $L_n^{\unsure{\flat|\natural}}$ le sont.

e) Si $k$ est noeth. régulier, $L_\bullet$ est $\partial$-parfait sss les $L_n$ sont de t. f. et \deleted{\unsure{la} $\unsure{\ldots}$ \unsure{fn.}} $\forall\,\unsure{k|K}$ corps résiduel, la fonction
% archive sheet number: 13

\scanpage{15}
\[
\unsure{n}\longmapsto\mathrm{rg}_K L_n\otimes_k K
\]
est polynômiale \deleted{pour $n$ grand}. [le degré du polynôme \inserted{(supposé non nul, i.e. $L_\bullet\ne0$)} est alors le plus grand entier $\unsure{+n|n|N}$ tel que $L_n^!\ne0$, et les coefficients du polynôme exprimé comme combinaison linéaire des $P_i(n)=\binom n i$ $(i\geq0)$ sont les $\mathrm{rg}_K L_i^!\otimes_k K$ ; les coefficients dudit polynôme (par r. aux $P_i$, ou pour l'écriture habituelle \deleted{et} ou comme comb. lin. des \unsure{$n\mapsto n^i$}) sont des fonctions constructibles en le pt variable $\unsure{\mathfrak p|p}\in\operatorname{Spec} k$].

\uline{NB} Si dans e) on suppose $k$ noeth mais non régulier, la condition \inserted{\unsure{énoncée}} de $\partial$-perfection est tjrs néc, mais pas suffisante -- il faut une condition néc. et suff. plus subtile, savoir l'existence d'un $L'_\bullet\unsure{\longrightarrow|\mathrel{-}}L_\bullet$, avec $L'_\bullet$ \uline{str.} $\partial$-parfait (cf c)) et que \deleted{soit un \unsure{isom.}} \inserted{\unsure{qu-i|quis}} i.e. induise un iso sur les $\pi_i$. (De plus, \unsure{$\exists\exists$|$\exists$} $L'\to L$ épi)

% slanted
\margin{Stabilité des notions \unsure{préc.} pour sommes, \unsure{extensions} ; \unsure{et pour} produits \unsure{tensoriels} $\Lambda^i,\Gamma^i,\operatorname{Sym}^i$ ! dans le cas \uuline{str.} $\partial$ ps. coh. et \uuline{str.} $\partial$ parfaits \deleted{\unsure{c'est évident}}}

On veut un théorème : \uline{\unsure{Pour|Dévis.}} \uline{\unsure{tout|précis}} des $L_\bullet$ \unsure{str. $\partial$-ps. coh.} \deleted{\unsure{on}} -- cela n'est possible bien sûr que \deleted{plus subtilt de \unsure{ps.} $\partial$-ps. coh.}
% archive sheet number: 14

\scanpage{16}
si on dispose d'une théorie analogue pour les complexes de chaînes \deleted{\unsure{de proj. de}} ps. coh, ce qui n'est guère le cas que si on suppose $k$ principal, (Dedekind à la rigueur). Supposons donc $k$ \deleted{\unsure{noeth.}} principal \deleted{\illegible{2}} \inserted{(le cas $k$ un corps est déjà traité, avec la notion des $L_\bullet$ spéciaux et str. spéciaux)}
\inserted{donc \unsure{bien entendu} \unsure{d'un|qu'un} projectif de t.f. est libre}
, et soit $\operatorname{Max}(k)$ l'ens des idéaux \deleted{premiers} \inserted{maximaux} de $k$, $\unsure{\mathscr D|\mathfrak D|D|\mathfrak B}=\mathbb N^{(\operatorname{Max} k)}-\{0\}$

\uline{Th.} Tout $L_\bullet$ \inserted{str.} $\partial$-ps. coh. est isomorphe à un produit de modules de la forme suivante

a) $\deleted{M_i\otimes_k}(\Lambda^i\Phi_\bullet)^{\mu_i}\quad(i\geq1)$ \deleted{$M_i$ $k$-module libre \deleted{libre} de t.f.} \inserted{$\mu_i\in\mathbb N$}

b) $\deleted{N_j\otimes_k}(\Lambda^j\Psi_*)^{\nu_j}\quad(j\geq0)\quad\deleted{N_j}\quad(\nu_j\in\mathbb N)$

c) $\deleted{\unsure{P_{i,\ell}}\otimes_k}\Phi(i,\unsure{\ell|b})_*^{\unsure{f_{i,\ell}|\beta_{i,\ell}|h_{i,\ell}|\rho_{i,\ell}}}\quad(i\geq1)\quad\deleted{\unsure{P_{i,\ell}}}\quad f_{i,\ell}\in\mathbb N,$

\deleted{où $\Phi(i,\ell)_{\unsure{*}}$}

$\ell$ un \deleted{idéal} diviseur de $k$ (cf plus bas la définition de $\Phi(i,\ell)_*$) ; de façon précise
\[
\begin{aligned}
L_\bullet\simeq{}&\prod_{i\geq1}M_i\otimes_k\Lambda^i\Phi_\bullet
 \times\prod_{j\geq0}N_j\otimes_k\Lambda^j\Psi_*\\
&{}\times\prod_{\substack{i\geq1\\\ell\in\mathscr D}}\Phi_{i,\ell}^{f_{i,\ell}}
\end{aligned}
\]
où pour tout $i\geq1$, \deleted{\unsure{l'ens des div.}} les $f_{i,\ell}$ sauf un nb fini sont nuls. (NB de tels produits sont bien str. $\partial$-ps coh.) On peut supposer de plus que pour $i$ donné, l'ens des $\ell\in\mathscr D$ \deleted{\unsure{non}} tels que $f_{i,\ell}\ne0$ soit tlt ordonné, moyennant
% archive sheet number: 15
